\section{后续影响与扩展}

ResNet的思想深刻影响了后续深度学习架构的发展，催生了一系列重要变体：

\begin{itemize}
    \item \textbf{ResNeXt}：引入分组卷积（Group Convolution）来增加残差块的宽度，在不显著增加参数量下提升性能，成为后来EfficientNet等工作的基础。
    \item \textbf{DenseNet}：将残差连接发展为稠密连接（Dense Connection），每一层都直接与所有后续层连接，进一步缓解了梯度消失并促进了特征复用。
    \item \textbf{Wide ResNet}：在保持残差结构的同时增加通道数（宽度）而减少深度，表明宽度在一定条件下比深度更为重要。
    \item \textbf{ResNet变体}：如Pre-activation ResNet（BN和ReLU前置）、ResNet with Squeeze-and-Excitation（SE模块）等，在残差块中引入注意力机制，进一步提升了表示能力。
\end{itemize}

此外，残差学习框架已超越图像识别领域，广泛应用于自然语言处理（如Transformer中的跳跃连接）、语音识别、强化学习（如ResNet用于Atari游戏），甚至图神经网络中。几乎任何需要堆叠多层的深度学习模型都可以受益于残差连接。

ResNet的发展也标志着深度学习从“手工设计特征”转向“手工设计网络结构”时代的成熟，后续的DenseNet、SENet、EfficientNet等无不是在残差范式上的创新。可以说，ResNet不仅是一个具体模型，更是一种设计哲学——它证明了只要解决优化问题，深度就能带来收益。